% !TEX program = pdflatex
\documentclass[10pt]{article}
\usepackage[utf8]{inputenc}\usepackage[T1]{fontenc}\usepackage[english]{babel}
\usepackage{lmodern,microtype}\usepackage[margin=.72in]{geometry}
\usepackage{amsmath,amssymb,amsthm,mathtools,booktabs,array,xcolor,enumitem}
\usepackage[most]{tcolorbox}\usepackage{fancyhdr}\usepackage[hidelinks]{hyperref}
\definecolor{CourseBlue}{HTML}{1A3C68}\definecolor{CourseGreen}{HTML}{216E39}
\definecolor{CourseGold}{HTML}{B8860B}\definecolor{CourseRed}{HTML}{9A2C2C}
\definecolor{SoftBlue}{HTML}{F2F6FA}\definecolor{SoftGreen}{HTML}{F1F8F3}
\definecolor{SoftGold}{HTML}{FFF8E8}\definecolor{SoftRed}{HTML}{FFF3F3}
\setlength{\parindent}{0pt}\setlength{\parskip}{.32em}
\setlist[itemize]{leftmargin=1.5em,itemsep=.12em,topsep=.2em}
\setlist[enumerate]{leftmargin=1.7em,itemsep=.12em,topsep=.2em}
\pagestyle{fancy}\fancyhf{}\lhead{Stochastic Calculus and Option Pricing}
\rhead{Lecture 12 Notes}\cfoot{\thepage}\renewcommand{\headrulewidth}{.3pt}
\newcommand{\E}{\mathbb E}\newcommand{\Q}{\mathbb Q}\newcommand{\cF}{\mathcal F}
\newcommand{\dd}{\,\mathrm d}\newcommand{\Var}{\operatorname{Var}}
\theoremstyle{definition}\newtheorem{definition}{Definition}[section]
\newtheorem{proposition}[definition]{Proposition}
\newtcolorbox{boardexample}[1]{enhanced,breakable,colback=SoftBlue,colframe=CourseBlue,boxrule=.55pt,arc=1mm,title={Example --- #1},fonttitle=\bfseries}
\newtcolorbox{studentexercise}[1]{enhanced,breakable,colback=SoftGold,colframe=CourseGold,boxrule=.55pt,arc=1mm,title={Exercise --- #1},fonttitle=\bfseries}
\newtcolorbox{shortsolution}[1][]{enhanced,breakable,colback=SoftGreen,colframe=CourseGreen,boxrule=.55pt,arc=1mm,title={Short solution#1},fonttitle=\bfseries}
\newtcolorbox{keyidea}[1]{enhanced,breakable,colback=white,colframe=CourseGreen,boxrule=.7pt,arc=1mm,title={#1},fonttitle=\bfseries}
\newtcolorbox{warningbox}[1]{enhanced,breakable,colback=SoftRed,colframe=CourseRed,boxrule=.65pt,arc=1mm,title={#1},fonttitle=\bfseries}
\title{\textbf{Beyond Black--Scholes}\\\large Lecture 12 Notes --- Local Volatility, Stochastic Volatility, and Stochastic Rates}
\author{}\date{}
\begin{document}\maketitle\vspace{-1.4em}
\begin{keyidea}{Goals and logical structure}
A useful model must explain the prices and risks that matter:
\[
\text{implied-volatility surface}\longrightarrow\text{local volatility}
\longrightarrow\text{random volatility}\longrightarrow\text{random discounting}.
\]
The trade-off is fit, realism, hedging, and tractability.
\end{keyidea}

\section{Why Black--Scholes is not enough}
Under Black--Scholes,
\[
\frac{\dd S_t}{S_t}=r\dd t+\sigma\dd W_t^\Q,
\qquad\sigma\text{ constant}.
\]
If this were exact, implied volatility would be constant across strikes and
maturities. Markets instead quote a surface
\((K,T)\mapsto\sigma_{\rm imp}(K,T)\), with smiles, skews, and term
structure.
\begin{boardexample}{equity skew}
For a fixed maturity, equity implied volatility typically satisfies
\[
\sigma_{\rm imp}(K_{\rm low})>
\sigma_{\rm imp}(K_{\rm ATM})>
\sigma_{\rm imp}(K_{\rm high}).
\]
Low-strike puts insure severe declines; their demand and price are high.
Implied volatility is the Black--Scholes parameter matching a quote, not a
literal historical volatility of that option.
\end{boardexample}
\begin{keyidea}{Three extensions, three questions}
\begin{center}
\begin{tabular}{p{.27\textwidth}|p{.28\textwidth}|p{.28\textwidth}}
\toprule
Observed issue & Extension & Main purpose\\\midrule
Static vanilla smile & local volatility & fit prices today\\
Volatility dynamics & stochastic volatility & model future smile and vol risk\\
Yield-curve risk & stochastic rates & model bonds and discounting\\
\bottomrule
\end{tabular}
\end{center}
\end{keyidea}

\section{Local volatility fits today's surface}
\subsection{One-factor state-dependent variance}
\begin{definition}[Local-volatility model]
Under \(\Q\),
\[
\frac{\dd S_t}{S_t}=r\dd t+
\sigma_{\rm loc}(t,S_t)\dd W_t.
\]
The instantaneous variance depends deterministically on time and current
spot. The model remains one-factor Markovian.
\end{definition}
For European payoff \(g(S_T)\), the price solves
\[
u_t+rsu_s+\frac12\sigma_{\rm loc}(t,s)^2s^2u_{ss}-ru=0,
\qquad u(T,s)=g(s).
\]
Only the constant variance of Black--Scholes has been replaced.

\subsection{Dupire calibration}
With constant rates and no dividends, an arbitrage-free call-price surface
\(C(T,K)\) gives formally
\[
\boxed{
\sigma_{\rm loc}(T,K)^2=
\frac{C_T(T,K)+rK C_K(T,K)}
{\frac12K^2C_{KK}(T,K)}.}
\]
Strike convexity is crucial:
\[
C_{KK}(T,K)=e^{-rT}p_{S_T}^{\Q}(K)\geq0.
\]
\begin{warningbox}{Differentiate a clean surface, not raw quotes}
Dupire uses a second strike derivative. Noise in a sparse quote table is
amplified and can give negative or unstable local variance. First build a
smooth, arbitrage-free price or implied-volatility surface.
\end{warningbox}
\begin{studentexercise}{local or extra state?}
Does local volatility introduce an additional random factor besides \(S_t\)?
\end{studentexercise}
\begin{shortsolution}
No. \(\sigma_{\rm loc}(t,S_t)\) is a deterministic function of the current
time and stock. The state remains \((t,S_t)\).
\end{shortsolution}
\begin{keyidea}{Strength and limitation}
Local volatility can reproduce a clean grid of vanilla prices today. It
also prescribes future smile motion through a deterministic spot-time rule;
that future dynamics need not be realistic for exotic hedging.
\end{keyidea}

\section{Stochastic volatility makes variance a risk factor}
\subsection{A two-factor model}
A generic specification is
\[
\frac{\dd S_t}{S_t}=r\dd t+\sqrt{V_t}\dd W_t^{(1)},\qquad
\dd V_t=b(t,V_t)\dd t+a(t,V_t)\dd W_t^{(2)},
\]
\[
\dd\langle W^{(1)},W^{(2)}\rangle_t=\rho\dd t.
\]
The Markov state is now \((S_t,V_t)\). Two paths with the same spot can have
different option values if their current variance levels differ.

\begin{boardexample}{Heston}
\[
\frac{\dd S_t}{S_t}=r\dd t+\sqrt{V_t}\dd W_t^{(1)},\qquad
\dd V_t=\kappa(\theta-V_t)\dd t+\eta\sqrt{V_t}\dd W_t^{(2)}.
\]
\(\theta\) is the long-run variance, \(\kappa\) mean-reversion speed,
\(\eta\) volatility of volatility, and \(\rho\) spot-variance correlation.
\end{boardexample}

The conditional mean illustrates mean reversion:
\[
\E[V_T\mid V_t=v]=\theta+(v-\theta)e^{-\kappa(T-t)}.
\]
In equities, \(\rho<0\) captures the leverage effect: a negative stock
shock tends to accompany higher variance, making downside puts relatively
expensive and generating skew.

\begin{warningbox}{Richer dynamics imply hedging risk}
There are two sources of randomness but often only stock and cash are
traded. The model is generally incomplete: stock delta alone cannot
perfectly hedge variance risk. Liquid vanilla options are often used as
additional risk-management instruments.
\end{warningbox}
\begin{studentexercise}{choose a volatility model}
Which model is more natural for a payoff sensitive to future realized
volatility: local or stochastic volatility?
\end{studentexercise}
\begin{shortsolution}
Stochastic volatility, because it retains a random variance state whose
future dynamics can affect the payoff and hedge.
\end{shortsolution}

\section{Stochastic rates make discounting random}
\subsection{Bonds and the yield curve}
With short rate \(r_t\),
\[
D(t,T)=\exp\left(-\int_t^T r_s\dd s\right),\qquad
P(t,T)=\E^\Q[D(t,T)\mid\cF_t].
\]
The curve can be parameterized equivalently by
\[
P(t,T)=e^{-R(t,T)(T-t)},\qquad
R(t,T)=-\frac{\log P(t,T)}{T-t},\qquad
f(t,T)=-\partial_T\log P(t,T).
\]
\begin{boardexample}{a curve calculation}
If \(P(0,1)=.96\) and \(P(0,2)=.90\), then
\[
R(0,2)=-\tfrac12\log(.90)=5.27\%,
\qquad L(0;1,2)=\frac{.96}{.90}-1=6.67\%.
\]
The second-year forward rate exceeds the two-year average zero rate.
\end{boardexample}

\subsection{Short-rate and forward-rate models}
A short-rate model specifies \(r_t\) and derives bond prices. Examples:
\[
\dd r_t=a(b-r_t)\dd t+\sigma\dd W_t\quad\text{(Vasicek)},
\]
\[
\dd r_t=a(b-r_t)\dd t+\sigma\sqrt{r_t}\dd W_t\quad\text{(CIR)}.
\]
HJM starts directly from the forward curve,
\[
\dd f(t,T)=\alpha(t,T)\dd t+\sigma(t,T)\dd W_t.
\]
Under the risk-neutral measure, no-arbitrage determines the drift
\(\alpha\) once forward-rate volatility is chosen.

\begin{studentexercise}{when are rates material?}
Is deterministic discounting a more defensible approximation for a
one-month vanilla equity option or for a five-year callable bond?
\end{studentexercise}
\begin{shortsolution}
The short-dated equity option is the more defensible case. A callable bond
is directly exposed to the future yield curve and an exercise decision, so
stochastic rates are material.
\end{shortsolution}

\section{Calibration and model choice}
\begin{keyidea}{Calibration is a market-pricing step}
Local volatility fits an arbitrage-free vanilla surface. Stochastic
volatility fits it through parameters such as
\(\kappa,\theta,\eta,\rho,V_0\). Rate models fit the initial curve together
with liquid caps, swaptions, or bond-option prices. Calibration aligns a
model with tradable prices before it is used for valuation or risk.
\end{keyidea}
\begin{warningbox}{Use the least complex model that captures material risk}
Local volatility is a natural first benchmark for accurate current vanilla
fit. A long-dated equity-linked payoff with rate exposure may need
stochastic volatility and stochastic rates. More factors mean more
calibration burden and more hedging uncertainty.
\end{warningbox}
\begin{studentexercise}{final model choice}
A desk must hedge a barrier option and reproduce the current vanilla smile
surface accurately. What is the natural first benchmark, and what caveat
remains?
\end{studentexercise}
\begin{shortsolution}
Local volatility is the natural first benchmark because it fits today's
vanilla surface. Its future smile dynamics can be unrealistic, so hedge
performance remains model dependent.
\end{shortsolution}
\begin{keyidea}{Final synthesis}
\[
\boxed{\text{Market prices}\longrightarrow\text{model assumptions}
\longrightarrow\text{calibration}\longrightarrow\text{pricing and hedging}.}
\]
A richer model earns its place only by explaining a material source of risk.
\end{keyidea}
\end{document}
